\NeedsTeXFormat{LaTeX2e}[2005/12/01]
%%    2009/03/12 v1.0 GAUBM Vorlage f�r Aschlussarbeiten Physik
%% Template fuer Bachelor- und Masterarbeiten
%% an der Fakultaet fuer Physik (c) Thomas Pruschke der GA Universit�t
%% Verbesserungsvorschlaege bitte an pruschke@theorie.physik.uni-goettingen.de
%%
%% Benoetigte Pakete: datenumber
%%

%%%%%%%%%%%%%%%%%%%%%%%%%%%%%%%%%%%%%%%%%%%%%%%%%%%%%%%%%%%%%%%%%%%%%%
%%%%%%%%%% Bitte vor dem Veraendern diese Datei umbenennen! %%%%%%%%%%
%%%%%%%%%%%%%%%%%%%%%%%%%%%%%%%%%%%%%%%%%%%%%%%%%%%%%%%%%%%%%%%%%%%%%%

%% scrbook - Ersatz f�r LaTeX book Klasse aus dem KOMA Script
%% Moegliche Optionen: diejenigen der Klasse scrbook ausser titlepage

%% deutsche Arbeit:
\documentclass[bachelor,       %% Typ der Arbeit: bachelor oder master
               twoside,        %% zweiseitiges Layout
               BCOR10mm,       %% Bindekorrektur 10 mm
%               liststotoc,nomtotoc,bibtotoc, %% Aufnahme der div. Verzeichnisse
                                              %% ins Inhaltsverzeichnis
%               english,ngerman, %% Alternativspr. Englisch, Dokumentspr. Deutsch
               ngerman,english  %% Alternativspr. Deutsch, Dokumentspr. Englisch
%               final,          %% Endversion; draft fuer schnelles Kompilieren
               ]{GAUBM}

\usepackage{setspace}  %% Zur Setzung des Zeilenabstandes
\usepackage{babel}     %% Sprachen-Unterstuetzung
\usepackage{calc}      %% ermoeglicht Rechnen mit Laengen und Zaehlern
\usepackage[T1]{fontenc}       %% Unterstutzung von Umlauten etc.
\usepackage[latin1]{inputenc}  %% 
%% in aktuellem Linux & MacOS X wird standardmaessig UTF8 kodiert!
%\usepackage[utf8]{inputenc}    %% Wenn latin1 nicht geht ...

\usepackage{amsmath,amssymb} %% zusaetzliche Mathe-Symbole

\usepackage{lmodern} %% type1-taugliche CM-Schrift als Variante zur
                     %% "normalen" EC-Schrift
%% Paket fuer bibtex-Datenbanken
\usepackage[comma,numbers,sort&compress]{natbib}
\bibliographystyle{plainnat}

\newcommand{\tabheadfont}[1]{\textbf{#1}} %% Tabellenkopf in Fett
\usepackage{booktabs}                      %% Befehle fuer besseres Tabellenlayout
\usepackage{longtable}                     %% umbrechbare Tabellen
\usepackage{array}                         %% zusaetzliche Spaltenoptionen

%% umfangreiche Pakete fuer Symbole wie \micro, \ohm, \degree, \celsius etc.
\usepackage{textcomp,gensymb}

%\usepackage{SIunits} %% Korrektes Setzen von Einheiten
\usepackage{units}   %% Variante fuer Einheiten

%% Hyperlinks im Dokument; muss als eines der letzten Pakete geladen werden
\usepackage[pdfstartview=FitH,      % Oeffnen mit fit width
            breaklinks=true,        % Umbrueche in Links, nur bei pdflatex default
            bookmarksopen=true,     % aufgeklappte Bookmarks
            bookmarksnumbered=true  % Kapitelnummerierung in bookmarks
            ]{hyperref}

%% Weiter benoetigte Pakete: datenumber
%% Falls dieses Paket nicht in der Installation vorhanden ist,
%% kann es von der Seite mit diesem Template heruntergeladen werden
%% und in einem LaTeX bekanntem Verzeichnis installiert werden (notfalls
%% dem Verzeichnis mit der Arbeit).

%%%%%%%%%%%%%% EIGENE VER�NDERUNGEN:#####################
%% Bra-Ket Notation: \bracket{||} oder \bra{}
\usepackage{braket}

\usepackage{cleveref}
\usepackage{bm}
\usepackage{algpseudocode}
\usepackage{algorithmicx}

\newcommand{\vect}[1]{\bm{#1}}
\renewcommand{\d}[1]{\ensuremath{\operatorname{d}\!{#1}}}
%%########################################################


\begin{document}
%%
%%                   Ab hier muessen die Anpassungen geschehen
%%
%% Hier den eigenen Namen einsetzen
\ThesisAuthor{Eric}{Bertok}
%% Hier den Geburtsort einsetzen
\PlaceOfBirth{Kassel}
%% Titel Arbeit. Das erste Argument ist der deutsche, das zweite der
%% englische Titel.
\ThesisTitle{Hier steht das Thema der Arbeit in deutsch}{Here comes the title of the thesis in english}
%% Erst- und Zweitgutacher/in
%% Ist der/die Betreuer/in nicht identisch mit dem/r Erstgutachter/in,
%% muss diese/r als optionales Argument angegeben werden.
\FirstReferee[Dr.\ \ldots]{Prof.\ Dr.\ Kehrein}
\Institute{Institut f\"ur Theoretische Physik}
\SecondReferee{Prof.\ Dr.\ Kree}
%% Beginn und Ende des Anfertigungszeitraumes
\ThesisBegin{1}{4}{2009}
\ThesisEnd{15}{7}{2009}
%% DO NOT TOUCH THESE LINES!!!!
\frontmatter
\maketitle
\cleardoublepage
%% Zusammenfassung. Falls nicht gewuenscht, bitte auskommentieren.
\begin{abstract}
  Here the key results of the thesis can be presented in about
  half a page.
  \bigskip\par
  \textbf{Keywords:} Physics, Bachelor thesis
  \end{abstract}
%% So laesst sich in die andere Sprache umschalten (Englisch bzw. Deutsch)
\begin{otherlanguage}{ngerman}
\begin{abstract}
  Hier werden auf einer halben Seite die Kernaussagen der Arbeit
  zusammengefasst.
%% Optional: Stichwoerter. Wenn nicht gewuenscht, koennen die beiden
%% folgenden Zeilen geloescht werden
  \bigskip\par
  \textbf{Stichw�rter:} Physik, Bachelorarbeit

\end{abstract}
\end{otherlanguage}

%% Ende des Vorspanns
\cleardoublepage
%% Ab hier 1 1/2 facher Zeilenabstand (durch setspace-Paket)
\onehalfspacing
%% Erzeugt Inhaltsverzeichnis
\tableofcontents

%% Hier kann man seine Bezeichnungsweisen erklaeren. Falls nicht
%% benoetigt, bis einschliesslich \end{nomenclature} auskommentieren
\begin{nomenclature}
%% Fuer die Berechnung der Spaltenbreiten muss \usepackage{calc}
%% geladen sein!
\section*{Lateinische Buchstaben}
\noindent
\begin{longtable}[l]{p{0.2\textwidth}p{0.7\textwidth-6\tabcolsep}p{0.1\textwidth}}
  \tabheadfont{Variable}&\tabheadfont{Bedeutung}&\tabheadfont{Einheit}\\\midrule\endhead
  $A$ & Querschnittsfl"ache & $\unit{m^2}$\\
  $c$ & Geschwindigkeit & $\unitfrac{m}{s}$
\end{longtable}
\section*{Griechische Buchstaben}
\begin{longtable}[l]{p{0.2\textwidth}p{0.7\textwidth-6\tabcolsep}p{0.1\textwidth}}
  \tabheadfont{Variable}&\tabheadfont{Bedeutung}&\tabheadfont{Einheit}\\\midrule\endhead
  $\alpha$  & Winkel & $\unit{\degree}$; --\\
  $\varrho$ & Dichte & $\unitfrac{kg}{m^3}$
\end{longtable}
\section*{Indizes}
\begin{longtable}[l]{p{0.2\textwidth}p{0.8\textwidth-4\tabcolsep}}
  \tabheadfont{Index}&\tabheadfont{Bedeutung}\\\midrule\endhead
  m & Meridian\\
  $r$ & Radial
\end{longtable}
\section*{Abk"urzungen}
\begin{longtable}[l]{p{0.2\textwidth}p{0.8\textwidth-4\tabcolsep}}
  \tabheadfont{Abk"urzung}&\tabheadfont{Bedeutung}\\\midrule\endhead
  2D & zweidimensional\\
  3D & dreidimensional\\
  max & maximal
\end{longtable}
\end{nomenclature}
%% \listoftables und \listoffigures sollten nur bei genuegender Anzahl Tabellen
%% verwendet werden
%\listoffigures
%\listoftables

\mainmatter   %% Anfang Hauptteil

\chapter{Einleitung}
Diese Vorlage \verb!GAUBM!
f�r Bachelor- bzw.\ Masterarbeiten ist eine �berarbeitung der
Vorlage von Simon Dreher f�r Abschlu�arbeiten am
Institut f�r Mikrosystemtechnologie (IMTEK)
an der Universit�t Freiburg. Die eigentliche Datei mit der Klassendefinition
ist \verb!GAUBN.cls!, die Sie zusammen mit dieser Datei erhalten haben. Weitere
Dateien sind \verb!datenumber.sty! und die zugeh�rigen Sprachdefinitionen
\verb!\datenumber*.ldf!. Im Verzeichnis \verb!figures! finden sich die
von der Klasse ben�tigten Logos (Universit�t und Physik) sowie Beispielbilder
f�r die �bersetzung dieser Beispieldatei (\verb!bthesis.tex!).
Sie k�nnen diese Datei als Vorlage f�r Ihre Arbeit nutzen und entsprechend
modifizieren. Bitte denken Sie daran, sie vorher unter einem eigenen Namen
abzuspeichern.
Um die Datei anzupassen, gehen Sie wie folgt vor:

Bei den Parametern zu \verb!\documentclass[...]{GAUBM}! in der Pr�ambel
kann man durch Umschalten
zwischen \verb!english,ngerman! und \verb!ngerman,english! eine
deutsche Arbeit (erste Variante) mit Englisch als Alternativsprache bzw.\ eine
englische Arbeit (zweite Variante) mit deutsch als Alternativsprache
w�hlen. Im laufenden Text kann man mit 
\begin{verbatim}
\begin{otherlanguage}{english/ngerman}
...
\end{otherlanguage}
\end{verbatim}
zur alternativen Sprache wechseln.

Nach \verb!\begin{document}! m�ssen zuerst ein paar Befehle mit
Information �ber die Arbeit aufgerufen werden:
\begin{enumerate}
\item \verb!\ThesisAuthor{Vorname}{Nachname}!: Die Argumente sind der
Vorname und Nachname der Autorin bzw.\ des Autors der Arbeit.
\item \verb!\PlaceOfBirth{Wohnort}!: Der Geburtsort der Autorin bzw.\ des Autors.
\item \verb!\ThesisTitle{Deutscher Titel}{English title}!: Der deutsche und englische Titel der Arbeit gem�� Antrag.
\item \verb!\Institute{Institut}!: Das Institut, an dem die Arbeit angefertigt wurde.
\item \verb!\FirstReferee[Betreuer/in]{Erste/r Gutachter/in}!: Voller Titel
und Name des/r Erstgutachter/in. Ist der Betreuer der Arbeit \emph{nicht}
identisch mit dem/r Erstgutachter/in, so mu� der volle Titel und der Name
des/r Betreuer/in als optionales Argument in eckigen Klammern erscheinen.
\item \verb!\SecondReferee{Zweite/r Gutachter/in}!: Voller Titel
und Name des/r Zweitgutachter/in.
\item \verb!\ThesisBegin{Tag}{Monat}{Jahr}!: Datum des Beginns der Anfertigung
der Arbeit gem�� Antrag.
\item \verb!\ThesisEnd{Tag}{Monat}{Jahr}!: Datum der Fertigstellung der Arbeit.
\item Optional kann mit
\begin{verbatim} 
\begin{abstract}
...
\end{abstract}
\end{verbatim}
eine maximal eine halbe Seite lange Zusammenfassung eingef�gt werden.

Falls man die Zusammenfassung in der alternativen Sprache verfassen m�chte,
dann geht das mit der Befehlsfolge
\begin{verbatim} 
\begin{otherlanguage}{english/ngerman}
\begin{abstract}
...
\end{abstract}
\end{otherlanguage}
\end{verbatim}
\end{enumerate}

\chapter{Grundlagen}
In diesem Kapitel werden die theoretischen Grundlagen erl�utert.

Wichtige Gleichungen, die in der Arbeit h�ufiger zitiert werden,
sollten eine Gleichungsnummer erhalten.
\begin{equation}
  \label{eq:pythagoras}
  a^2+b^2=c^2
\end{equation}
Zum Beispiel wird in Gleichung~\ref{eq:pythagoras} der Satz des Pythagoras
angegeben.

Gerade im Bereich der Grundlagen wird viel Literatur zitiert, z.B.\
\cite{Menz97}. Falls
mehrere Literaturzitate auf einmal zitiert werden, ist folgendes
z.B.\ m�glich \cite{Horn90,DINEN6232,Menz97,Knuth84}.

\section{Unterkapitel Gliederungsebene 2}
Hier sollte etwas Text stehen.
\subsection{Unterkapitel Gliederungsebene 3}
Noch ein paar Beispiele zu Abbildungen und Tabellen:

Abbildung~\ref{fig:bildplatzhalter} verdeutlicht \dots

Wie die Abb.~\ref{fig:bildplatzhalter} und
Tab.~\ref{tab:tabellenplatzhalter} verdeutlichen \dots

\begin{figure}
  \centering
  \includegraphics[width=0.5\linewidth]{figures/bild}
  \caption{Bildbeschreibung}
  \label{fig:bildplatzhalter}
\end{figure}

Text\dots
\begin{table}
  \centering
  \begin{tabular}{llll}
    \toprule
    $A$-Wert&$B$-Wert&$C$-Wert&$D$-Wert\\
    \midrule
    aaaaaa&bbbbbbb&cccccc&ddddddd\\
    aaaaaa&bbbbbbb&cccccc&ddddddd\\
    \bottomrule
  \end{tabular}
  \caption{Tabellenbeschreibung}
  \label{tab:tabellenplatzhalter}
\end{table}

Text\dots

\chapter{``Lab course''}
\section{Theoretical foundation and numerical methods}
\subsection{Introduction, exact diagonalization}
In what follows, we are interested in a one-dimensional XXZ Heisenberg chain with next-nearest neighbour coupling. As such, we consider the following Hamiltonian:
\begin{align}
    H(J,\mu)&=H_1(J,\mu)+\lambda H_2(J,\mu),\\
    H_n(J,\mu)&=J \sum_{i=1}^{N}\left[\frac{1}{2}\left( S_i^+ S_{i+n}^- + S_i^- S_{i+n}^+ \right)+\mu S_i^z S_{i+n}^z\right],
    \label{eq:XXZham}
\end{align}
where the spins are labeled $i=1,\dots,N$ and $J$ is the interaction constant. The scalar product for the spin operators $\mathbf{S}_i\cdot\mathbf{S}_{i+n}$ has been decomposed into creation-annihilation operators $S^\pm=S^x \pm i S^y$ while  $\mu$ introduces an anisotropy for the z-direction compared to the x and y directions.
Finally, $\lambda$ is a parameter controlling the strength of the next-nearest neighbor (NNN) interactions cf. \cite[4ff.]{xxz_springer}.\\
First, would like to numerically compute the time evolution of such a system by means of exact diagonalization.
\subsection{Block diagonalization}
In the following, techniques and notation from \cite[55ff.]{sand11} are used.
The states are denoted as $\ket{S_0^z,S_1^z,\dots,S_{N-1}}$, where the subscript $i=0,\dots,N-1$ refers to the site i of our one-dimensional chain. A particular state is abridgedly labeled by arrows, e.g.
$\ket{\uparrow\downarrow\uparrow\downarrow\dots}$.
Furthermore, periodic boundary conditions, which are useful for the construction of symmetries but not needed, are assumed ($i=N-1+n\equiv n-1\label{eq:periodic}$).
Exact diagonalization allows for total knowledge of the system by choosing a basis, setting up a hamiltonian and diagonalizing it numerically. While theoretically applicable to any system size, this method is limited to around 30 spins, as the computational cost for diagonalization is $\propto 8^\text{dim. of basis}$[citation needed.] and, without utilizing symmetries, the spin basis above grows like $2^N$.\\
Symmetries can be used to transform the hamiltonian matrix into block-diagonal form cf. \cite[56]{sand11} with each block belonging to states with a conserved quantum number, thereby reducing the basis dimension. In this first chapter, only the conservation of z-component of the total spin ($\hat{=}$ total magnetization) $m^z=\sum_{i=0}^{N-1} S_i^z$ is used:\\
It is easy to verify, using the normal commutation relations for the spin operators [citation needed] that
\begin{align}
    [S_i^z,H]&=\frac{1}{2}(A_1+A_2),\\
    A_n&=S_i^+ S_{i+n}^- -S_i^- S_{i+n}^+ -S_{i-n}^+ S_i^- +S_{i-n}^- S_i^+.
    \label{eq:comm}
\end{align}
By using the periodic boundary conditions it follows that
\begin{align}
    \sum_{i=0}^{N-1}A_n=\sum_{i=0}^{N-1}\left( S_i^+ S_{i+n}^- -S_i^- S_{i+n}^+ \right)-\underbrace{\sum_{i=n}^{N-1+n}}_{\equiv\sum_{i=0}^{N-1}}\left( S_i^+ S_{i+n}^- -S_i^- S_{i+n}^+ \right)=0,
    \label{eq:Ancomm}
\end{align}
and therefore
\begin{align}
    \left[H,m^z\right]=0,
    \label{eq:totalmagcomm}
\end{align}
meaning that $m^z$ is a conserved quantity. By constructing a basis with fixed magnetization, blocks with the same $S^z$ can be diagonalized independently, meaning that $S^z$ is a conserved quantity. By constructing a basis with fixed magnetization, blocks with the same $S^z$ can be considered independently.

\subsection{Bit representation}
As Spin-$1/2$ partlicles only have two possible states $(\uparrow \& \downarrow)$, they can be elegantly represented in a computer via the bits ``1'' and ``0''.
More precisely, A lattice state $\ket{S_0^z,S_1^z,\dots,S_{N-1}}$ is represented as the bit representation of an unsigned integer, with the first bit denoted by the index 0. \cite[59]{sand11}, e.g. $\ket{\uparrow\downarrow\downarrow\uparrow}\hat{=} ~1001~ \dot{=}~9$. Symmetries that will come in later, like spin translation or rotation manifest themselves as operations on these bit representations.\\

\subsection{Construction of the fixed magnetization basis}
A Basis list of states having the same magnetization $m_z$ is found by iterating over all integers $s$ from the $2^n$ dimensional basis and counting the amount of ``1'' bits $n_\uparrow$, checking whether they amount to the target magnetization $m_z$=$n_\uparrow-1/2$.\\
With $s[i]$ labeling the $i$'th bit, starting from 0, the following pseudocode is used to obtain an ordered list of basis states $s_a$:

\begin{algorithmic}
    \State $a=0$
    \ForAll{$s=0..2^N-1$}
    \If{$\sum_i s[i]=n_\uparrow$}
        \State$ a \gets a+1$
        \State$s_a \gets s$
        \EndIf
    \EndFor
    \State $M\gets a$
\end{algorithmic}\label{pseudo:basismag}
$M$ is the basis dimension. Thus, a natural basis is given by all vectors 
\begin{align}
    \ket{a}=\hat{e}_a=(0,\dots,0,\underbrace{1}_{a},0,\dots,0)^T, a=1..M,     \label{eq:basismag}
\end{align}
each representing the corresponding integer state $s_a$.

\subsection{Construction of the hamilton matrix}
The hamiltonian (\ref{eq:XXZham}) is translated into a matrix for the fixed magnetization basis cf. \cite[60f.]{sand11}.
The diagonal terms are $\braket{a|S_i^z S_{i+n}^z|a}=\pm1/4~\mu$ ,since the $i$'th and $i+n$'th spin are either equal or different.
The off-diagonal terms $\left( S_i^+ S_{i+n}^- + S_i^- S_{i+n}^+ \right)$ flip the bits $i$ and $i+n$ $(0\leftrightarrow1)$ and amount to zero otherwise.
A bisectional search [Pseudocode?] is used to find position $b$ of the flipped state $s*$, giving a matrix element contribution of $\braket{b|( S_i^+ S_{i+n}^- + S_i^- S_{i+n}^+ )|a}=+1/2$.
In pseudocode, the matrix belonging to the NNN hamiltonian ($n=1,2$) is constructed as follows:
\begin{algorithmic}
    \For{$a=1..M$}
    \For{$n=1,2$}
    \State $j=(i+n)\mod N$
    \State \If{$a[i]=a[j]$}
    \State $H(a,a)\gets H(a,a)+\frac{1}{4}~\mu$
    \Else
    \State $H(a,a)\gets H(a,a)-\frac{1}{4}~\mu$
    \State $s*=\text{flipstate}(a)$
    \State $b\gets$ findstate($s*$)
    \State $H(a,b)=\frac{1}{2}$
    \EndIf
    \EndFor
    \EndFor
\end{algorithmic}

\subsection{Diagonalization and time evolution}
The Armadillo package \cite{Armadillo} is used for matrix diagonalization and vector operations.







\chapter{Ergebnisse}
Test\dots
citation: \cite{sand11}
F5 - Test:\\
\begin{align}
    \bra{\Psi}\ket{\Phi}
    \label{H}
\end{align}
in \cref{H} wird\dots
$\braket{8|s|7|6}$\\

\textbf{Griechische Buchstaben:}\\
\begin{align}
    \alpha\beta\chi\delta\varepsilon\varphi\gamma\eta\iota\kappa\lambda\mu\nu\pi\theta\rho\sigma\tau\upsilon\varsigma\xi\psi\zeta\\
    \Delta\Phi\Gamma\Lambda\Xi\Psi\Sigma\Upsilon\Omega\label{1}\\
    \braket{\Psi|\hat{H}|\Psi}\\\label{2}\\
    \vect{V}=\vec{V}\\
\end{align}
Label \cref{1,2}

$(\partial_\mu\partial^\mu+m^2)\Psi=0$\\
$\frac{5}{8}=z$\\
$A\circ B$\\
$A\equiv B$
\begin{align}
    \setminus\bigcap\times\wedge\bigcap\bigcup\subset\supset\le\ge\nonumber\dot{A}\ddot{B}\sqrt{A}\Big|\int_{-\infty}^{\infty}\d x  
    \label{b}
\end{align}

\begin{align}
    \partial_\mu F^{\mu\nu}&=\mu_0 j^\nu\\
    \partial_\mu \tilde{F}^{\mu\nu}&=0
\end{align}
\chapter{Diskussion}
Text\dots
\section{Unterkapitel}
Text\dots
\subsection{Unterkapitel}
Text\dots

\chapter{Zusammenfassung}
Text\dots

\appendix
\chapter{Calculation of $[S_i^z,H]=0$}

\chapter{zweiter Anhang}
Text\dots

\cleardoublepage
%% Bibliographie. Das Argument muss der Name der BIBTeX-Datenbank stehen.
%% Ein Beispiel fuer eine solche Datenbank finden Sie in bthesis_datenbank.bib
\bibliography{bthesis_datenbank} 

\chapter*{Danksagung}
Dank\dots

%% Dieser Befehl MUSS am Ende stehen und erzeugt die Erklaerung ueber die
%% benutzten Mittel
\Declaration
\end{document}
